% =============================================================================
% Chapitre 3 : Réalisation
% =============================================================================

\chapter{Réalisation}
\label{chap:realisation}

\section*{Introduction}
Ce chapitre constitue le volet technique principal de notre travail. Il présente en détail
la mise en \oe uvre concrète du robot de tests automatisés pour la plateforme SpeedLine.
Nous y décrivons l'environnement de développement retenu, puis nous détaillons
l'implémentation de chaque couche de la pyramide de tests : tests unitaires, tests API,
tests de bout en bout (E2E), et tests de performance. Nous abordons ensuite les aspects
transversaux que sont les tests de sécurité, le \textit{seed} de la base de données,
l'intégration CI/CD et le reporting avec Allure.

% =============================================================================
\section{Environnement de développement}
\label{sec:env-dev}
% =============================================================================

La mise en place du robot de tests nécessite un ensemble d'outils couvrant les différents
niveaux de la pyramide de tests. Le tableau~\ref{tab:env-dev} récapitule les principaux
outils utilisés, leur version ainsi que leur rôle au sein du projet.

\begin{table}[H]
\centering
\caption{Environnement de développement et outils utilisés}
\label{tab:env-dev}
\begin{tabularx}{\textwidth}{l l X}
\toprule
\textbf{Outil} & \textbf{Version} & \textbf{Rôle} \\
\midrule
Node.js          & 24.x   & Environnement d'exécution JavaScript pour les tests Playwright et k6 \\
npm              & 11.x   & Gestionnaire de paquets Node.js \\
Java (OpenJDK)   & 17     & Langage du backend Spring Boot (microservices) \\
Maven            & 3.9    & Outil de build et de gestion des dépendances Java \\
Playwright       & 1.49   & Framework de tests API et E2E pour les applications web \\
JUnit 5          & 5.x    & Framework de tests unitaires Java \\
Mockito          & 5.x    & Bibliothèque de mocking pour les tests unitaires Java \\
Jasmine + Karma  & ---    & Frameworks de tests unitaires pour les frontends Angular \\
Patrol           & 3.13   & Framework de tests d'intégration pour les applications Flutter \\
k6               & 1.7    & Outil de tests de charge et de performance \\
Allure           & 2.27   & Générateur de rapports de tests interactifs \\
Docker           & ---    & Conteneurisation des services pour l'environnement de test \\
Git + GitLab CI  & ---    & Gestion de versions et intégration continue \\
VS Code          & ---    & Éditeur de code pour TypeScript, JavaScript et Dart \\
IntelliJ IDEA    & ---    & IDE pour le développement Java (backend) \\
Flutter          & 3.41   & Framework de développement des applications mobiles \\
Angular          & 17     & Framework de développement des interfaces web (admin-panel, partner-dashboard) \\
\bottomrule
\end{tabularx}
\end{table}

Le choix de ces outils s'inscrit dans une logique de cohérence avec la pile technologique
existante du projet SpeedLine. Les tests unitaires backend reposent sur JUnit~5 et Mockito,
standards de l'écosystème Spring Boot. Pour les tests API et E2E web, Playwright a été
retenu pour sa rapidité d'exécution, sa fiabilité et son support multi-navigateurs.
Les applications mobiles Flutter sont testées avec Patrol, qui offre la possibilité
d'interagir avec les dialogues natifs du système d'exploitation. Enfin, k6 a été choisi
pour les tests de performance grâce à sa syntaxe JavaScript familière et ses métriques
intégrées.

% =============================================================================
\section{Mise en place des tests unitaires (Couche 1)}
\label{sec:tests-unitaires}
% =============================================================================

Les tests unitaires constituent la base de notre pyramide de tests. Ils visent à vérifier
le comportement isolé de chaque composant logiciel, en simulant les dépendances externes
à l'aide de \textit{mocks}. Cette approche garantit des tests rapides, déterministes et
indépendants de l'infrastructure.

% -----------------------------------------------
\subsection{Tests unitaires backend (JUnit 5 + Mockito)}
\label{subsec:tests-unitaires-backend}
% -----------------------------------------------

\subsubsection{Approche adoptée}

Pour chaque service du backend, nous appliquons la même méthodologie :

\begin{enumerate}
  \item \textbf{Isolation des dépendances} : les repositories, clients Feign et services
        externes sont simulés (\textit{mockés}) avec Mockito via l'annotation
        \texttt{@Mock}.
  \item \textbf{Injection automatique} : le service sous test est instancié avec
        \texttt{@InjectMocks}, ce qui injecte automatiquement les mocks.
  \item \textbf{Organisation en classes imbriquées} : les tests sont groupés par
        méthode testée à l'aide de l'annotation \texttt{@Nested}, améliorant
        la lisibilité des rapports.
  \item \textbf{Assertions expressives} : nous utilisons AssertJ pour des assertions
        fluides et lisibles.
\end{enumerate}

\subsubsection{Exemple : AuthServiceImplTest}

Le listing~\ref{lst:auth-service-test} illustre la structure d'un test unitaire pour
le service d'authentification. Ce test vérifie le scénario d'inscription réussie
d'un nouveau client.

\begin{lstlisting}[language=Java, caption={Test d'inscription -- AuthServiceImplTest.java}, label={lst:auth-service-test}]
@ExtendWith(MockitoExtension.class)
class AuthServiceImplTest {

    @Mock private UserRepository userRepo;
    @Mock private PasswordEncoder passwordEncoder;
    @Mock private JwtTokenProvider tokenProvider;
    @Mock private AuthenticationManager authenticationManager;
    @Mock private UserServiceClient userServiceClient;

    @InjectMocks
    private AuthServiceImpl authService;

    @Nested
    @DisplayName("register()")
    class RegisterTests {

        @Test
        @DisplayName("Should register a new customer successfully")
        void register_success_customer() {
            RegisterRequest request = new RegisterRequest();
            request.setEmail("new@example.com");
            request.setPassword("Password123");
            request.setFirstName("Jane");
            request.setLastName("Doe");
            request.setPhoneNumber("+21655443322");
            request.setRole(Role.CUSTOMER);

            when(userRepo.existsByEmailIgnoreCase("new@example.com"))
                .thenReturn(false);
            when(passwordEncoder.encode("Password123"))
                .thenReturn("encoded-pass");
            when(userRepo.save(any(User.class)))
                .thenAnswer(invocation -> {
                    User u = invocation.getArgument(0);
                    u.setId(10L);
                    return u;
                });

            authService.register(request);

            ArgumentCaptor<User> captor =
                ArgumentCaptor.forClass(User.class);
            verify(userRepo).save(captor.capture());
            User saved = captor.getValue();

            assertThat(saved.getEmail())
                .isEqualTo("new@example.com");
            assertThat(saved.getPassword())
                .isEqualTo("encoded-pass");
            assertThat(saved.getRole())
                .isEqualTo(Role.CUSTOMER);
            assertThat(saved.getStatus())
                .isEqualTo(UserStatus.PENDING);
            verify(userServiceClient).createCustomer(any());
        }
    }
}
\end{lstlisting}

On note l'utilisation de \texttt{ArgumentCaptor} pour capturer l'objet \texttt{User}
passé à la méthode \texttt{save()}, permettant de vérifier finement les valeurs
assignées lors de l'inscription. La méthode \texttt{verify()} s'assure que l'appel
au service distant \texttt{userServiceClient} a bien été effectué.

\subsubsection{Exemple : JwtServiceTest}

Le listing~\ref{lst:jwt-service-test} présente les tests du fournisseur de jetons JWT.
Contrairement aux tests précédents, ces tests n'utilisent pas de mocks car le
\texttt{JwtTokenProvider} est testé avec sa logique réelle de génération et de
validation de jetons.

\begin{lstlisting}[language=Java, caption={Test de génération JWT -- JwtServiceTest.java}, label={lst:jwt-service-test}]
class JwtServiceTest {

    private JwtTokenProvider jwtTokenProvider;
    private static final String TEST_SECRET =
        "my-super-secret-key-for-testing-jwt-tokens-"
        + "must-be-long-enough-256-bits";
    private static final long EXPIRATION_MS = 3600000L;

    private User testUser;

    @BeforeEach
    void setUp() {
        jwtTokenProvider = new JwtTokenProvider();
        ReflectionTestUtils.setField(
            jwtTokenProvider, "jwtSecret", TEST_SECRET);
        ReflectionTestUtils.setField(
            jwtTokenProvider, "jwtExpirationMs", EXPIRATION_MS);

        testUser = User.builder()
                .id(42L)
                .email("test@example.com")
                .role(Role.CUSTOMER)
                .status(UserStatus.ACTIVE)
                .build();
    }

    @Nested
    @DisplayName("generateToken()")
    class GenerateTokenTests {

        @Test
        @DisplayName("Should embed email as subject in token")
        void generateToken_containsEmailAsSubject() {
            String token = jwtTokenProvider.generateToken(testUser);
            String email = jwtTokenProvider
                .getEmailFromToken(token);
            assertThat(email).isEqualTo("test@example.com");
        }

        @Test
        @DisplayName("Should embed role claim in token")
        void generateToken_containsRole() {
            String token = jwtTokenProvider.generateToken(testUser);
            String role = jwtTokenProvider
                .getRoleFromToken(token);
            assertThat(role).isEqualTo("CUSTOMER");
        }
    }
}
\end{lstlisting}

\subsubsection{Bilan des tests unitaires backend}

Le tableau~\ref{tab:unit-tests-backend} présente la répartition des tests unitaires
par microservice. Au total, \textbf{419 tests} ont été implémentés, couvrant l'ensemble
des services métier de la plateforme.

\begin{table}[H]
\centering
\caption{Répartition des tests unitaires backend par microservice}
\label{tab:unit-tests-backend}
\begin{tabular}{l c l}
\toprule
\textbf{Microservice} & \textbf{Nombre de tests} & \textbf{Classes testées} \\
\midrule
auth-service         & 100 & AuthServiceImpl, JwtService, AuthController \\
order-service        &  87 & OrderService, CartService, PriceCalculation \\
user-service         &  57 & CustomerService, CourierService \\
delivery-service     &  41 & DispatchCycleService, EligibilityFilter, Solvers \\
payment-service      &  30 & PaymentServiceImpl, WalletService, RefundService \\
location-service     &  29 & ZoneServiceImpl, GeolocationServiceImpl \\
promotion-service    &  20 & PromotionApplicationService, Validators \\
notification-service &  18 & NotificationServiceImpl \\
analytics-service    &  16 & AnalyticsServiceImpl, ReportServiceImpl \\
support-service      &  11 & SupportTicketServiceImpl \\
review-service       &  10 & ReviewServiceImpl \\
\midrule
\textbf{Total}       & \textbf{419} & \\
\bottomrule
\end{tabular}
\end{table}

Les services les plus critiques --- \texttt{auth-service} et \texttt{order-service} ---
concentrent le plus grand nombre de tests (100 et 87 respectivement), reflétant leur
complexité fonctionnelle et leur importance dans le flux métier de la plateforme.

% -----------------------------------------------
\subsection{Tests unitaires frontend (Jasmine + Karma)}
\label{subsec:tests-unitaires-frontend}
% -----------------------------------------------

Les applications web Angular (\texttt{admin-panel} et \texttt{partner-dashboard})
disposent de tests unitaires exécutés par Karma avec le framework d'assertions Jasmine.
Ces tests vérifient le comportement des services, des gardes de route (\textit{guards})
et des intercepteurs HTTP.

\subsubsection{Exemple : AuthService (admin-panel)}

Le listing~\ref{lst:auth-service-spec} montre un extrait des tests du service
d'authentification du panneau d'administration. Le service \texttt{ApiService} est simulé
avec \texttt{jasmine.createSpyObj} pour isoler la logique d'authentification.

\begin{lstlisting}[language=Java, caption={Test du service d'authentification Angular -- auth.service.spec.ts}, label={lst:auth-service-spec}]
describe('AuthService', () => {
  let service: AuthService;
  let apiServiceSpy: jasmine.SpyObj<ApiService>;

  const mockLoginResponse = {
    access_token: 'eyJhbGciOiJIUzI1NiJ9...',
    user: {
      id: '1',
      email: 'admin@speedline-test.com',
      role: 'SUPER_ADMIN',
    },
  };

  beforeEach(() => {
    apiServiceSpy = jasmine.createSpyObj('ApiService',
      ['get', 'post', 'put']);
    TestBed.configureTestingModule({
      imports: [HttpClientTestingModule],
      providers: [
        AuthService,
        { provide: ApiService, useValue: apiServiceSpy },
      ],
    });
    service = TestBed.inject(AuthService);
    localStorage.clear();
    sessionStorage.clear();
  });

  describe('login', () => {
    it('should store tokens on successful login', (done) => {
      apiServiceSpy.post.and.returnValue(
        of(mockLoginResponse));
      service.login({
        email: 'admin@test.com',
        password: 'pass',
        rememberMe: false
      }).subscribe({
        next: () => {
          expect(service.isLoggedIn()).toBeTrue();
          expect(service.currentUser()?.email)
            .toBe('admin@speedline-test.com');
          done();
        },
      });
    });

    it('should reject non-admin roles', (done) => {
      const customerResponse = {
        ...mockLoginResponse,
        user: { ...mockLoginResponse.user,
                role: 'CUSTOMER' },
      };
      apiServiceSpy.post.and.returnValue(
        of(customerResponse));
      service.login({
        email: 'customer@test.com',
        password: 'pass',
        rememberMe: false
      }).subscribe({
        error: (err) => {
          expect(err.message).toBe('ACCESS_DENIED');
          done();
        },
      });
    });
  });
});
\end{lstlisting}

\subsubsection{Exemple : AuthGuard}

Le listing~\ref{lst:auth-guard-spec} présente les tests du garde de route
\texttt{authGuard}, qui protège les pages nécessitant une authentification.

\begin{lstlisting}[language=Java, caption={Test du garde d'authentification -- auth.guard.spec.ts}, label={lst:auth-guard-spec}]
describe('authGuard', () => {
  let authServiceSpy: jasmine.SpyObj<AuthService>;
  let routerSpy: jasmine.SpyObj<Router>;

  beforeEach(() => {
    authServiceSpy = jasmine.createSpyObj('AuthService',
      ['isAuthenticated']);
    routerSpy = jasmine.createSpyObj('Router', ['navigate']);
    TestBed.configureTestingModule({
      providers: [
        { provide: AuthService, useValue: authServiceSpy },
        { provide: Router, useValue: routerSpy },
      ],
    });
  });

  it('should allow access when authenticated', () => {
    authServiceSpy.isAuthenticated.and.returnValue(true);
    const result = TestBed.runInInjectionContext(
      () => authGuard(mockRoute, mockState));
    expect(result).toBeTrue();
  });

  it('should redirect to login when not authenticated',
    () => {
    authServiceSpy.isAuthenticated.and.returnValue(false);
    const result = TestBed.runInInjectionContext(
      () => authGuard(mockRoute, mockState));
    expect(result).toBeFalse();
    expect(routerSpy.navigate).toHaveBeenCalledWith(
      ['/auth/login'],
      { queryParams: { returnUrl: '/dashboard' } });
  });
});
\end{lstlisting}

\subsubsection{Exemple : AuthInterceptor}

Le listing~\ref{lst:auth-interceptor-spec} montre les tests de l'intercepteur HTTP,
qui ajoute automatiquement le jeton JWT aux requêtes sortantes.

\begin{lstlisting}[language=Java, caption={Test de l'intercepteur HTTP -- auth.interceptor.spec.ts}, label={lst:auth-interceptor-spec}]
describe('authInterceptor', () => {
  let httpClient: HttpClient;
  let httpMock: HttpTestingController;
  let authServiceSpy: jasmine.SpyObj<AuthService>;

  beforeEach(() => {
    authServiceSpy = jasmine.createSpyObj('AuthService',
      ['getToken']);
    TestBed.configureTestingModule({
      providers: [
        provideHttpClient(
          withInterceptors([authInterceptor])),
        provideHttpClientTesting(),
        { provide: AuthService,
          useValue: authServiceSpy },
      ],
    });
    httpClient = TestBed.inject(HttpClient);
    httpMock = TestBed.inject(HttpTestingController);
  });

  it('should add Authorization header', () => {
    authServiceSpy.getToken.and.returnValue('my-jwt');
    httpClient.get('/api/v1/orders').subscribe();
    const req = httpMock.expectOne('/api/v1/orders');
    expect(req.request.headers.get('Authorization'))
      .toBe('Bearer my-jwt');
  });

  it('should skip auth endpoints', () => {
    authServiceSpy.getToken.and.returnValue('my-jwt');
    httpClient.get('/api/v1/auth/login').subscribe();
    const req = httpMock.expectOne('/api/v1/auth/login');
    expect(req.request.headers.has('Authorization'))
      .toBeFalse();
  });
});
\end{lstlisting}

\subsubsection{Bilan des tests unitaires frontend}

Le tableau~\ref{tab:unit-tests-frontend} présente le résultat des tests unitaires
frontend ainsi que la couverture de code mesurée par Istanbul via Karma.

\begin{table}[H]
\centering
\caption{Résultats des tests unitaires frontend}
\label{tab:unit-tests-frontend}
\begin{tabular}{l c c}
\toprule
\textbf{Application} & \textbf{Nombre de tests} & \textbf{Couverture de code} \\
\midrule
admin-panel         & 26 & 57,78\% \\
partner-dashboard   &  8 & 34,82\% \\
\midrule
\textbf{Total}      & \textbf{34} & --- \\
\bottomrule
\end{tabular}
\end{table}

La couverture de 57,78\% pour l'\texttt{admin-panel} couvre les services critiques
(authentification, RBAC, intercepteurs). Le \texttt{partner-dashboard}, avec 34,82\%,
se concentre sur les flux d'authentification essentiels. Ces taux constituent une base
solide, avec un potentiel d'amélioration lors des itérations futures.

% =============================================================================
\section{Mise en place des tests API (Couche 2)}
\label{sec:tests-api}
% =============================================================================

Les tests API constituent la deuxième couche de notre pyramide. Ils vérifient le
comportement des microservices à travers leurs points d'entrée REST, en envoyant des
requêtes HTTP réelles à l'\textit{API Gateway} et en validant les réponses.

% -----------------------------------------------
\subsection{Architecture des tests API}
\label{subsec:archi-tests-api}
% -----------------------------------------------

Pour faciliter l'écriture et la maintenance des tests API, nous avons développé une
classe utilitaire \texttt{ApiClient} qui encapsule la logique d'authentification et de
requêtage. Le listing~\ref{lst:api-client} présente cette classe.

\begin{lstlisting}[language=Java, caption={Classe utilitaire ApiClient -- api-client.ts}, label={lst:api-client}]
import { APIRequestContext, request } from '@playwright/test';

export class ApiClient {
  private baseUrl: string;
  private token: string | null = null;

  constructor(baseUrl?: string) {
    this.baseUrl = baseUrl
      || process.env.API_GATEWAY_URL
      || 'http://localhost:8080';
  }

  async login(email: string, password: string)
    : Promise<string> {
    const ctx = await request.newContext(
      { baseURL: this.baseUrl });
    const response = await ctx.post('/api/auth/login', {
      data: { email, password },
    });
    if (!response.ok()) {
      throw new Error(
        `Login failed: ${response.status()}`);
    }
    const body = await response.json();
    this.token = body.accessToken || body.token;
    await ctx.dispose();
    return this.token!;
  }

  async authenticatedContext()
    : Promise<APIRequestContext> {
    if (!this.token) {
      throw new Error('Not logged in.');
    }
    return request.newContext({
      baseURL: this.baseUrl,
      extraHTTPHeaders: {
        Authorization: `Bearer ${this.token}`,
      },
    });
  }

  async get(path: string) {
    const ctx = await this.authenticatedContext();
    const res = await ctx.get(path);
    await ctx.dispose();
    return res;
  }

  async post(path: string, data: unknown) {
    const ctx = await this.authenticatedContext();
    const res = await ctx.post(path, { data });
    await ctx.dispose();
    return res;
  }
}
\end{lstlisting}

Cette classe offre plusieurs avantages :
\begin{itemize}
  \item \textbf{Centralisation de l'authentification} : un seul appel à
        \texttt{login()} suffit pour obtenir un jeton réutilisable.
  \item \textbf{Gestion automatique des en-têtes} : le jeton JWT est
        automatiquement ajouté à chaque requête via
        \texttt{authenticatedContext()}.
  \item \textbf{Nettoyage des ressources} : chaque contexte Playwright est
        libéré après utilisation avec \texttt{dispose()}.
\end{itemize}

% -----------------------------------------------
\subsection{Tests par service}
\label{subsec:tests-par-service}
% -----------------------------------------------

\subsubsection{Tests du service d'authentification}

Le listing~\ref{lst:auth-api-spec} présente les tests API du service
d'authentification. Ces tests vérifient les scénarios de connexion, de validation
de jeton et de gestion des erreurs.

\begin{lstlisting}[language=Java, caption={Tests API du service d'authentification -- auth.api.spec.ts}, label={lst:auth-api-spec}]
test.describe('API - Auth Service', () => {
  test('POST /api/auth/login - valid credentials',
    async ({ request }) => {
    const response = await request.post(
      `${API_BASE}/api/auth/login`, {
      data: {
        email: TEST_ACCOUNTS.admin.email,
        password: TEST_ACCOUNTS.admin.password,
      },
    });
    expect(response.status()).toBe(200);
    const body = await response.json();
    expect(body.accessToken || body.token).toBeTruthy();
  });

  test('POST /api/auth/login - invalid credentials',
    async ({ request }) => {
    const response = await request.post(
      `${API_BASE}/api/auth/login`, {
      data: {
        email: 'nonexistent@test.com',
        password: 'WrongPassword123!',
      },
    });
    expect(response.status())
      .toBeGreaterThanOrEqual(400);
    expect(response.status()).toBeLessThan(500);
  });

  test('GET /api/auth/me - user info with valid token',
    async ({ request }) => {
    const loginRes = await request.post(
      `${API_BASE}/api/auth/login`, {
      data: {
        email: TEST_ACCOUNTS.admin.email,
        password: TEST_ACCOUNTS.admin.password,
      },
    });
    const { accessToken } = await loginRes.json();
    const meRes = await request.get(
      `${API_BASE}/api/auth/me`, {
      headers: {
        Authorization: `Bearer ${accessToken}` },
    });
    expect(meRes.status()).toBe(200);
  });
});
\end{lstlisting}

\subsubsection{Tests du service de commandes}

Le listing~\ref{lst:orders-api-spec} illustre les tests CRUD du service de commandes,
incluant la pagination et la création de commandes.

\begin{lstlisting}[language=Java, caption={Tests API du service de commandes -- orders.api.spec.ts}, label={lst:orders-api-spec}]
test.describe('API - Order Service', () => {
  let api: ApiClient;
  let customerApi: ApiClient;

  test.beforeAll(async () => {
    api = new ApiClient();
    await api.login(TEST_ACCOUNTS.admin.email,
      TEST_ACCOUNTS.admin.password);
    customerApi = new ApiClient();
    await customerApi.login(TEST_ACCOUNTS.customer.email,
      TEST_ACCOUNTS.customer.password);
  });

  test('GET /api/orders - admin should list orders',
    async () => {
    const response = await api.get('/api/orders');
    expect(response.status()).toBe(200);
    const body = await response.json();
    expect(Array.isArray(body.content || body))
      .toBeTruthy();
  });

  test('POST /api/orders - customer should create order',
    async () => {
    const response = await customerApi.post(
      '/api/orders', {
      partnerId: 'test-partner-id',
      items: TEST_ORDER.items,
      deliveryAddress: TEST_ORDER.deliveryAddress,
      paymentMethod: TEST_ORDER.paymentMethod,
    });
    expect([200, 201, 400])
      .toContain(response.status());
  });
});
\end{lstlisting}

\subsubsection{Test du cycle de vie complet d'une commande}

Le listing~\ref{lst:e2e-order-flow} présente le test E2E API le plus complet : le cycle
de vie intégral d'une commande, de la création à la livraison. Ce test implique quatre
acteurs distincts (client, partenaire, coursier, administrateur) et valide les
transitions d'état.

\begin{lstlisting}[language=Java, caption={Test du cycle de vie complet -- e2e-order-flow.api.spec.ts}, label={lst:e2e-order-flow}]
test.describe('API - E2E Order Lifecycle', () => {
  let adminApi: ApiClient;
  let customerApi: ApiClient;
  let partnerApi: ApiClient;
  let courierApi: ApiClient;

  test.beforeAll(async () => {
    adminApi = new ApiClient();
    await adminApi.login(TEST_ACCOUNTS.admin.email,
      TEST_ACCOUNTS.admin.password);
    customerApi = new ApiClient();
    await customerApi.login(TEST_ACCOUNTS.customer.email,
      TEST_ACCOUNTS.customer.password);
    partnerApi = new ApiClient();
    await partnerApi.login(TEST_ACCOUNTS.partner.email,
      TEST_ACCOUNTS.partner.password);
    courierApi = new ApiClient();
    await courierApi.login(TEST_ACCOUNTS.courier.email,
      TEST_ACCOUNTS.courier.password);
  });

  test('full order lifecycle: create -> accept '
    + '-> pickup -> deliver', async () => {
    // 1. Get available partners
    const partnersRes = await customerApi.get(
      '/api/partners?status=ACTIVE');
    const partners = await partnersRes.json();
    const partnerId = partners.content[0].id;

    // 2. Create order
    const createRes = await customerApi.post(
      '/api/orders', {
      partnerId,
      items: TEST_ORDER.items,
      deliveryAddress: TEST_ORDER.deliveryAddress,
      paymentMethod: TEST_ORDER.paymentMethod,
    });
    const order = await createRes.json();

    // 3. Partner accepts order
    // 4. Courier picks up
    // 5. Courier delivers
    // 6. Order completed
  });
});
\end{lstlisting}

\subsubsection{Bilan des tests API}

Le tableau~\ref{tab:api-tests} récapitule l'ensemble des fichiers de tests API et le
nombre de cas de test par catégorie. Au total, \textbf{104 tests API} couvrent les
17 fichiers de spécifications.

\begin{table}[H]
\centering
\caption{Répartition des tests API par fichier de spécification}
\label{tab:api-tests}
\begin{tabular}{l c}
\toprule
\textbf{Fichier de test} & \textbf{Nombre de tests} \\
\midrule
auth.api.spec.ts            &  8 \\
orders.api.spec.ts          &  6 \\
users.api.spec.ts           &  6 \\
support.api.spec.ts         &  6 \\
partners.api.spec.ts        &  5 \\
delivery.api.spec.ts        &  5 \\
notifications.api.spec.ts   &  5 \\
reviews.api.spec.ts         &  5 \\
analytics.api.spec.ts       &  5 \\
payments.api.spec.ts        &  4 \\
location.api.spec.ts        &  4 \\
promotions.api.spec.ts      &  3 \\
e2e-order-flow.api.spec.ts  &  2 \\
contracts.api.spec.ts       &  9 \\
smoke.api.spec.ts           &  8 \\
security.api.spec.ts        & 16 \\
websocket.api.spec.ts       &  8 \\
\midrule
\textbf{Total (17 fichiers)} & \textbf{104} \\
\bottomrule
\end{tabular}
\end{table}

Les tests de sécurité représentent la catégorie la plus fournie avec 16 tests,
reflétant l'importance accordée à la sécurisation des API. Les tests
\texttt{contracts} vérifient la conformité des contrats d'API (formats de réponse,
codes HTTP), tandis que les tests \texttt{smoke} servent de vérification rapide
de la disponibilité des services.

% =============================================================================
\section{Mise en place des tests E2E (Couche 3)}
\label{sec:tests-e2e}
% =============================================================================

Les tests de bout en bout (\textit{End-to-End}) constituent la troisième couche de
la pyramide. Ils simulent le comportement réel des utilisateurs en interagissant avec
les interfaces graphiques web et mobiles.

% -----------------------------------------------
\subsection{Tests Web (Playwright)}
\label{subsec:tests-web}
% -----------------------------------------------

\subsubsection{Le patron Page Object Model}

Pour structurer les tests E2E web, nous adoptons le patron \textit{Page Object Model}
(POM). Chaque page de l'application est représentée par une classe TypeScript qui
encapsule les localisateurs d'éléments et les actions utilisateur. Le
listing~\ref{lst:admin-login-page} illustre cette approche pour la page de connexion
de l'administration.

\begin{lstlisting}[language=Java, caption={Page Object -- AdminLoginPage}, label={lst:admin-login-page}]
import { Page, Locator, expect } from '@playwright/test';

export class AdminLoginPage {
  readonly page: Page;
  readonly emailInput: Locator;
  readonly passwordInput: Locator;
  readonly loginButton: Locator;
  readonly errorMessage: Locator;

  constructor(page: Page) {
    this.page = page;
    this.emailInput = page.getByLabel(/email/i);
    this.passwordInput = page.getByLabel(/password/i);
    this.loginButton = page.getByRole('button',
      { name: /sign in|login|connexion/i });
    this.errorMessage = page.locator(
      '[class*="error"], [class*="alert-danger"]');
  }

  async goto() {
    await this.page.goto('/auth/login');
  }

  async login(email: string, password: string) {
    await this.emailInput.fill(email);
    await this.passwordInput.fill(password);
    await this.loginButton.click();
  }

  async expectDashboardRedirect() {
    await this.page.waitForURL('**/dashboard',
      { timeout: 15_000 });
    await expect(this.page).toHaveURL(/dashboard/);
  }

  async expectError() {
    await expect(this.errorMessage.first())
      .toBeVisible({ timeout: 5_000 });
  }
}
\end{lstlisting}

Les avantages de cette approche sont multiples :
\begin{itemize}
  \item \textbf{Maintenabilité} : si un sélecteur change, une seule classe doit
        être modifiée.
  \item \textbf{Réutilisabilité} : les actions (connexion, navigation) sont
        réutilisées dans plusieurs tests.
  \item \textbf{Lisibilité} : les tests décrivent des intentions utilisateur
        plutôt que des interactions DOM de bas niveau.
\end{itemize}

\subsubsection{Tests d'authentification (admin-panel)}

Le listing~\ref{lst:auth-e2e-spec} présente les tests E2E d'authentification de
l'administration, utilisant le \texttt{AdminLoginPage} défini précédemment.

\begin{lstlisting}[language=Java, caption={Tests E2E d'authentification -- auth.spec.ts}, label={lst:auth-e2e-spec}]
test.describe('Admin Panel - Authentication', () => {
  test.use({ storageState:
    { cookies: [], origins: [] } });

  test('should display login page', async ({ page }) => {
    const loginPage = new AdminLoginPage(page);
    await loginPage.goto();
    await expect(loginPage.emailInput).toBeVisible();
    await expect(loginPage.passwordInput).toBeVisible();
    await expect(loginPage.loginButton).toBeVisible();
  });

  test('should login with valid admin credentials',
    async ({ page }) => {
    const loginPage = new AdminLoginPage(page);
    await loginPage.goto();
    await loginPage.login(
      TEST_ACCOUNTS.admin.email,
      TEST_ACCOUNTS.admin.password);
    await loginPage.expectDashboardRedirect();
  });

  test('should show error with invalid credentials',
    async ({ page }) => {
    const loginPage = new AdminLoginPage(page);
    await loginPage.goto();
    await loginPage.login(
      'invalid@test.com', 'WrongPassword123!');
    await loginPage.expectError();
  });

  test('should redirect unauthenticated users',
    async ({ page }) => {
    await page.goto('/dashboard');
    await expect(page).toHaveURL(/auth/);
  });
});
\end{lstlisting}

\subsubsection{Tests du tableau de bord}

Le listing~\ref{lst:dashboard-e2e-spec} illustre les tests de navigation du tableau
de bord d'administration, vérifiant le chargement des pages et la navigation latérale.

\begin{lstlisting}[language=Java, caption={Tests E2E du tableau de bord -- dashboard.spec.ts}, label={lst:dashboard-e2e-spec}]
test.describe('Admin Panel - Dashboard', () => {
  test('should load dashboard after login',
    async ({ page }) => {
    const dashboard = new AdminDashboardPage(page);
    await dashboard.goto();
    await dashboard.expectLoaded();
  });

  test('should display sidebar navigation',
    async ({ page }) => {
    const dashboard = new AdminDashboardPage(page);
    await dashboard.goto();
    await expect(dashboard.sidebar).toBeVisible();
  });

  test('should navigate to orders',
    async ({ page }) => {
    const dashboard = new AdminDashboardPage(page);
    await dashboard.goto();
    await dashboard.navigateTo('orders');
    await expect(page).toHaveURL(/orders/);
  });

  test('should navigate to partners',
    async ({ page }) => {
    const dashboard = new AdminDashboardPage(page);
    await dashboard.goto();
    await dashboard.navigateTo('partners');
    await expect(page).toHaveURL(/partners/);
  });
});
\end{lstlisting}

\subsubsection{Tests de régression visuelle et d'accessibilité}

En complément des tests fonctionnels, nous avons mis en place :

\begin{itemize}
  \item \textbf{Tests de régression visuelle} : des captures d'écran de référence
        sont prises lors de la première exécution. Les exécutions suivantes comparent
        les captures actuelles avec les références pour détecter toute modification
        visuelle non intentionnelle. Ces tests sont implémentés dans les fichiers
        \texttt{visual-regression.spec.ts} pour chaque application web.
  \item \textbf{Tests d'accessibilité} : nous utilisons la bibliothèque
        \texttt{axe-core} intégrée à Playwright pour vérifier la conformité WCAG
        des pages. Les tests sont définis dans les fichiers
        \texttt{accessibility.spec.ts} et détectent les problèmes de contraste,
        de structure sémantique et d'attributs ARIA.
\end{itemize}

% -----------------------------------------------
\subsection{Tests Mobile (Flutter + Patrol)}
\label{subsec:tests-mobile}
% -----------------------------------------------

Les applications mobiles SpeedLine --- l'application client (\texttt{customer\_app})
et l'application coursier (\texttt{courier\_app}) --- sont développées avec Flutter.
Les tests d'intégration utilisent le framework \texttt{integration\_test} de Flutter,
complété par Patrol pour la gestion des dialogues natifs.

\subsubsection{Tests de l'application client}

Le listing~\ref{lst:customer-app-test} présente un extrait des tests d'intégration
de l'application client. Ces tests couvrent l'authentification, la navigation et
la gestion du panier.

\begin{lstlisting}[language=Java, caption={Tests d'intégration -- customer\_app/app\_test.dart}, label={lst:customer-app-test}]
void main() {
  IntegrationTestWidgetsFlutterBinding
    .ensureInitialized();

  group('Authentication', () {
    testWidgets('should display login screen',
      (tester) async {
      app.main();
      await tester.pumpAndSettle(
        const Duration(seconds: 3));
      final textFields = find.byType(TextFormField);
      expect(textFields, findsWidgets);
    });

    testWidgets('should login with valid credentials',
      (tester) async {
      app.main();
      await tester.pumpAndSettle(
        const Duration(seconds: 3));
      final textFields = find.byType(TextFormField);
      await tester.enterText(
        textFields.at(0), testConfig.customerEmail);
      await tester.enterText(
        textFields.at(1), testConfig.customerPassword);
      await tester.pumpAndSettle();
      final loginButtons =
        find.byType(ElevatedButton);
      await tester.tap(loginButtons.first);
      await tester.pumpAndSettle(
        const Duration(seconds: 5));
    });
  });

  group('Home & Browse', () => {
    testWidgets('should show bottom navigation bar',
      (tester) async {
      app.main();
      await _loginAsCustomer(tester);
      final bottomNav =
        find.byType(BottomNavigationBar);
      final navBar = find.byType(NavigationBar);
      expect(
        bottomNav.evaluate().isNotEmpty
        || navBar.evaluate().isNotEmpty,
        isTrue);
    });
  });
}
\end{lstlisting}

\subsubsection{Tests de l'application coursier}

Le listing~\ref{lst:courier-app-test} montre les tests de l'application coursier,
qui vérifient notamment le basculement en ligne/hors ligne.

\begin{lstlisting}[language=Java, caption={Tests d'intégration -- courier\_app/app\_test.dart}, label={lst:courier-app-test}]
void main() {
  IntegrationTestWidgetsFlutterBinding
    .ensureInitialized();

  group('Authentication', () {
    testWidgets('should login with courier credentials',
      (tester) async {
      app.main();
      await tester.pumpAndSettle(
        const Duration(seconds: 3));
      final textFields = find.byType(TextFormField);
      await tester.enterText(
        textFields.at(0), testConfig.courierEmail);
      await tester.enterText(
        textFields.at(1), testConfig.courierPassword);
      await tester.pumpAndSettle();
      final loginButtons =
        find.byType(ElevatedButton);
      await tester.tap(loginButtons.first);
      await tester.pumpAndSettle(
        const Duration(seconds: 5));
    });
  });

  group('Home Screen', () {
    testWidgets('should show online/offline toggle',
      (tester) async {
      app.main();
      await _loginAsCourier(tester);
      final toggle = find.byType(Switch);
      final switchWidget =
        find.byType(SwitchListTile);
      expect(
        toggle.evaluate().isNotEmpty
        || switchWidget.evaluate().isNotEmpty,
        isTrue);
    });

    testWidgets('should toggle courier availability',
      (tester) async {
      app.main();
      await _loginAsCourier(tester);
      final toggle = find.byType(Switch);
      if (toggle.evaluate().isNotEmpty) {
        await tester.tap(toggle.first);
        await tester.pumpAndSettle(
          const Duration(seconds: 2));
      }
    });
  });
}
\end{lstlisting}

\subsubsection{Gestion des dialogues natifs avec Patrol}

Patrol (version 3.13) est utilisé pour gérer les dialogues natifs du système
d'exploitation lors des tests mobiles. Il permet notamment de :

\begin{itemize}
  \item Accepter ou refuser les demandes de permission (localisation, notifications,
        caméra).
  \item Interagir avec les dialogues système (partage, sélecteur de fichiers).
  \item Gérer les alertes système qui ne sont pas accessibles via l'arbre de widgets
        Flutter standard.
\end{itemize}

% =============================================================================
\section{Mise en place des tests de performance (Couche 4)}
\label{sec:tests-performance}
% =============================================================================

Les tests de performance constituent la couche supérieure de notre pyramide de tests.
Ils évaluent le comportement du système sous différentes conditions de charge à l'aide
de l'outil k6.

% -----------------------------------------------
\subsection{Test de charge (\textit{Load Test})}
\label{subsec:test-charge}
% -----------------------------------------------

Le test de charge simule un scénario réaliste d'utilisation de la plateforme avec une
montée progressive du nombre d'utilisateurs virtuels. Le listing~\ref{lst:k6-load-test}
présente la configuration et les scénarios implémentés.

\begin{lstlisting}[language=Java, caption={Test de charge k6 -- k6-load-test.js}, label={lst:k6-load-test}]
import http from 'k6/http';
import { check, sleep, group } from 'k6';
import { Rate, Trend } from 'k6/metrics';

// Metriques personnalisees
const loginDuration = new Trend('login_duration');
const orderListDuration =
  new Trend('order_list_duration');
const errorRate = new Rate('errors');

export const options = {
  stages: [
    { duration: '30s', target: 10 },  // Montee
    { duration: '1m',  target: 25 },  // Montee
    { duration: '2m',  target: 25 },  // Plateau
    { duration: '1m',  target: 50 },  // Pic
    { duration: '1m',  target: 50 },  // Maintien pic
    { duration: '30s', target: 0 },   // Descente
  ],
  thresholds: {
    http_req_duration: ['p(95)<3000'],
    http_req_failed: ['rate<0.05'],
    login_duration: ['p(95)<2000'],
    errors: ['rate<0.1'],
  },
};

function login(email, password) {
  const res = http.post(
    `${BASE_URL}/api/auth/login`,
    JSON.stringify({ email, password }),
    { headers:
      { 'Content-Type': 'application/json' } }
  );
  loginDuration.add(res.timings.duration);
  const success = check(res, {
    'login status is 200': (r) => r.status === 200,
    'login has token': (r) => {
      const body = JSON.parse(r.body);
      return !!(body.accessToken || body.token);
    },
  });
  if (!success) { errorRate.add(1); return null; }
  errorRate.add(0);
  return JSON.parse(res.body).accessToken;
}

export default function () {
  const flow = Math.random();
  if (flow < 0.4)       adminFlow();
  else if (flow < 0.7)  customerBrowseFlow();
  else                   partnerFlow();
  sleep(1);
}
\end{lstlisting}

Le test de charge repose sur les principes suivants :

\begin{itemize}
  \item \textbf{Montée progressive} : le nombre d'utilisateurs virtuels augmente
        graduellement de 0 à 50, simulant une montée en charge réaliste.
  \item \textbf{Répartition des flux} : 40\% des utilisateurs exécutent le flux
        administrateur, 30\% le flux client et 30\% le flux partenaire.
  \item \textbf{Seuils de performance} : le test échoue si le 95e percentile
        dépasse 3 secondes ou si le taux d'erreur excède 5\%.
  \item \textbf{Métriques personnalisées} : des tendances dédiées mesurent les
        temps de réponse de la connexion et de la liste des commandes.
\end{itemize}

% -----------------------------------------------
\subsection{Test de stress (\textit{Stress Test})}
\label{subsec:test-stress}
% -----------------------------------------------

Le test de stress pousse le système au-delà de sa charge normale pour identifier les
points de rupture. Il monte progressivement jusqu'à 300 utilisateurs virtuels simultanés.

\begin{lstlisting}[language=Java, caption={Configuration du test de stress -- k6-stress-test.js}, label={lst:k6-stress-test}]
export const options = {
  stages: [
    { duration: '1m',  target: 50 },   // Charge normale
    { duration: '2m',  target: 100 },  // Charge haute
    { duration: '2m',  target: 200 },  // Charge stress
    { duration: '1m',  target: 300 },  // Point de rupture
    { duration: '2m',  target: 300 },  // Maintien
    { duration: '1m',  target: 0 },    // Recuperation
  ],
  thresholds: {
    http_req_duration: ['p(99)<5000'],
    errors: ['rate<0.3'],
  },
};
\end{lstlisting}

Le test de stress utilise \texttt{http.batch()} pour envoyer simultanément des requêtes
à plusieurs \textit{endpoints} (commandes, partenaires, utilisateurs), simulant ainsi
une utilisation intensive et concurrente du système. Le seuil d'erreur est relevé à 30\%
car l'objectif est d'observer le comportement de dégradation gracieuse (\textit{graceful
degradation}) plutôt que d'exiger une disponibilité totale.

% -----------------------------------------------
\subsection{Test de pic (\textit{Spike Test})}
\label{subsec:test-spike}
% -----------------------------------------------

Le test de pic simule un afflux soudain de trafic, tel qu'une heure de pointe ou une
promotion éclair. Il passe brutalement de 5 à 200 utilisateurs virtuels.

\begin{lstlisting}[language=Java, caption={Configuration du test de pic -- k6-spike-test.js}, label={lst:k6-spike-test}]
export const options = {
  stages: [
    { duration: '10s', target: 5 },    // Baseline
    { duration: '10s', target: 200 },  // Pic soudain
    { duration: '30s', target: 200 },  // Pic maintenu
    { duration: '10s', target: 5 },    // Fin du pic
    { duration: '1m',  target: 5 },    // Recuperation
    { duration: '10s', target: 0 },    // Arret
  ],
  thresholds: {
    http_req_duration: ['p(95)<5000'],
    errors: ['rate<0.2'],
  },
};
\end{lstlisting}

Ce test est particulièrement pertinent pour SpeedLine, car la plateforme de livraison
connaît des pics de trafic prévisibles (heures de repas) et imprévisibles (événements
promotionnels). La période de récupération d'une minute permet de vérifier que le
système retrouve son fonctionnement normal après le pic.

% =============================================================================
\section{Tests de sécurité}
\label{sec:tests-securite}
% =============================================================================

La sécurité étant un aspect critique d'une plateforme de livraison gérant des données
personnelles et des paiements, nous avons implémenté une suite de 16 tests de sécurité
couvrant les principales catégories de vulnérabilités.

\subsection{Injection SQL}

Les tests vérifient que les paramètres de recherche, les paramètres de chemin et les
corps de requête sont correctement assainis pour prévenir les injections SQL.

\begin{lstlisting}[language=Java, caption={Tests d'injection SQL -- security.api.spec.ts}, label={lst:sql-injection}]
test('should handle SQL injection in search params',
  async () => {
  const response = await adminApi.get(
    "/api/orders?search=' OR '1'='1");
  expect([200, 400]).toContain(response.status());
  if (response.status() === 200) {
    const body = await response.json();
    const items = body.content || body;
    expect(Array.isArray(items)).toBeTruthy();
  }
});

test('should handle SQL injection in POST body',
  async ({ request }) => {
  const loginRes = await request.post(
    `${API_BASE}/api/auth/login`, {
    data: {
      email: "admin@test.com' OR '1'='1",
      password: "' OR '1'='1",
    },
  });
  expect(loginRes.status())
    .toBeGreaterThanOrEqual(400);
});
\end{lstlisting}

\subsection{Cross-Site Scripting (XSS)}

Les tests XSS vérifient que les entrées utilisateur contenant du code JavaScript
malveillant sont correctement assainies ou rejetées par le serveur.

\begin{lstlisting}[language=Java, caption={Tests XSS -- security.api.spec.ts}, label={lst:xss-tests}]
test('should sanitize XSS in order creation',
  async () => {
  const response = await customerApi.post(
    '/api/orders', {
    partnerId: 'test-partner',
    items: [{ name: '<script>alert("xss")</script>',
              quantity: 1 }],
    deliveryAddress: {
      street: '<img src=x onerror=alert(1)>',
      city: 'Nabeul',
    },
    paymentMethod: 'CASH',
  });
  if (response.status() === 200
      || response.status() === 201) {
    const body = await response.text();
    expect(body).not.toContain('<script>');
  }
});
\end{lstlisting}

\subsection{Contournement d'authentification et contrôle d'accès}

Ces tests vérifient qu'un utilisateur ne peut pas accéder à des ressources non
autorisées, qu'un jeton invalide est correctement rejeté, et qu'un jeton est bien
invalidé après la déconnexion.

\begin{lstlisting}[language=Java, caption={Tests d'authentification et de contrôle d'accès}, label={lst:auth-bypass}]
test('should reject requests without auth token',
  async ({ request }) => {
  const response = await request.get(
    `${API_BASE}/api/orders`);
  expect(response.status()).toBe(401);
});

test('customer should not access admin endpoints',
  async () => {
  const response = await customerApi.get(
    '/api/users/admins');
  expect([401, 403]).toContain(response.status());
});

test('should invalidate token after logout',
  async ({ request }) => {
  // Login
  const loginRes = await request.post(
    `${API_BASE}/api/auth/login`, {
    data: { email: TEST_ACCOUNTS.admin.email,
            password: TEST_ACCOUNTS.admin.password },
  });
  const { accessToken } = await loginRes.json();

  // Logout
  await request.post(`${API_BASE}/api/auth/logout`, {
    headers: { Authorization: `Bearer ${accessToken}` },
  });

  // Old token should be blacklisted
  const afterLogout = await request.get(
    `${API_BASE}/api/auth/me`, {
    headers: { Authorization: `Bearer ${accessToken}` },
  });
  expect(afterLogout.status()).toBe(401);
});
\end{lstlisting}

\subsection{Limitation de débit et validation des entrées}

Les tests vérifient la présence d'un mécanisme de limitation de débit (\textit{rate
limiting}) sur le point d'entrée de connexion, ainsi que le rejet des charges utiles
surdimensionnées.

\begin{lstlisting}[language=Java, caption={Tests de limitation de débit}, label={lst:rate-limiting}]
test('should rate limit login attempts',
  async ({ request }) => {
  const attempts = [];
  for (let i = 0; i < 20; i++) {
    attempts.push(
      request.post(`${API_BASE}/api/auth/login`, {
        data: { email: 'brute-force@test.com',
                password: `attempt${i}` },
      })
    );
  }
  const responses = await Promise.all(attempts);
  const statusCodes = responses.map(r => r.status());
  const has429 = statusCodes.includes(429);
  if (!has429) {
    console.warn('WARNING: No rate limiting detected');
  }
});

test('should reject oversized request body',
  async ({ request }) => {
  const largePayload = 'x'.repeat(10 * 1024 * 1024);
  const response = await request.post(
    `${API_BASE}/api/auth/login`, {
    data: { email: largePayload, password: 'test' },
  });
  expect([400, 413, 414]).toContain(response.status());
});
\end{lstlisting}

% =============================================================================
\section{Seed de la base de données}
\label{sec:seed-database}
% =============================================================================

Avant l'exécution des tests, il est nécessaire de disposer d'un jeu de données
cohérent dans les différentes bases de données des microservices. Le script
\texttt{seed-test-data.js} automatise cette initialisation.

\subsection{Approche technique}

Le script se connecte successivement aux quatre bases de données principales
(\texttt{auth\_db}, \texttt{partner\_db}, \texttt{user\_db}, \texttt{location\_db})
et insère les données de test nécessaires. Les mots de passe sont hachés avec
l'algorithme \texttt{bcrypt} (facteur de coût 10) pour correspondre au format attendu
par le service d'authentification.

\begin{lstlisting}[language=Java, caption={Script de seed -- seed-test-data.js (extrait)}, label={lst:seed-script}]
const { Client } = require('pg');
const bcrypt = require('bcryptjs');

const TEST_USERS = [
  { email: 'admin@speedline-test.com',
    password: 'TestAdmin123!',
    firstName: 'Test', lastName: 'Admin',
    role: 'SUPER_ADMIN', phone: '+21611111111' },
  { email: 'partner@speedline-test.com',
    password: 'TestPartner123!',
    firstName: 'Test', lastName: 'Partner',
    role: 'PARTNER', phone: '+21622222222' },
  { email: 'customer@speedline-test.com',
    password: 'TestCustomer123!',
    firstName: 'Test', lastName: 'Customer',
    role: 'CUSTOMER', phone: '+21633333333' },
  { email: 'courier@speedline-test.com',
    password: 'TestCourier123!',
    firstName: 'Test', lastName: 'Courier',
    role: 'COURIER', phone: '+21644444444' },
];

async function seedAuthDB() {
  const client = new Client({
    ...DB_CONFIG, database: 'auth_db' });
  await client.connect();

  for (const user of TEST_USERS) {
    const passwordHash =
      await bcrypt.hash(user.password,
        await bcrypt.genSalt(10));
    await client.query(`
      INSERT INTO users (email, password, first_name,
        last_name, role, phone_number, status,
        is_email_verified, created_at, updated_at)
      VALUES ($1,$2,$3,$4,$5,$6,'ACTIVE',true,
        NOW(),NOW())
      ON CONFLICT (email) DO UPDATE SET
        password = $2, status = 'ACTIVE',
        is_email_verified = true
    `, [user.email, passwordHash, user.firstName,
        user.lastName, user.role, user.phone]);
  }
  await client.end();
}
\end{lstlisting}

\subsection{Comptes de test}

Le tableau~\ref{tab:test-accounts} récapitule les quatre comptes de test créés par le
script de seed. Ces comptes couvrent les quatre rôles principaux de la plateforme.

\begin{table}[H]
\centering
\caption{Comptes de test créés par le script de seed}
\label{tab:test-accounts}
\begin{tabular}{l l l}
\toprule
\textbf{Rôle} & \textbf{Email} & \textbf{Mot de passe} \\
\midrule
Administrateur  & admin@speedline-test.com    & TestAdmin123! \\
Partenaire      & partner@speedline-test.com  & TestPartner123! \\
Client          & customer@speedline-test.com & TestCustomer123! \\
Coursier        & courier@speedline-test.com  & TestCourier123! \\
\bottomrule
\end{tabular}
\end{table}

En plus des comptes utilisateur, le script crée un partenaire de test
(\og Test Restaurant Nabeul \fg{}) avec trois produits (Pizza Margherita, Coca Cola,
Salade Caesar), une zone de livraison (\og Zone Nabeul Centre \fg{}) et les profils
client et coursier correspondants dans la base \texttt{user\_db}. L'utilisation de
la clause \texttt{ON CONFLICT} garantit l'idempotence du script : il peut être exécuté
plusieurs fois sans créer de doublons.

% =============================================================================
\section{Intégration CI/CD}
\label{sec:integration-cicd}
% =============================================================================

L'intégration continue constitue un pilier essentiel de notre stratégie de tests.
Le fichier \texttt{.gitlab-ci-tests.yml} définit un ensemble de jobs qui s'intègrent
au pipeline existant du projet SpeedLine.

\subsection{Architecture du pipeline de tests}

Le pipeline est structuré en deux étapes (\textit{stages}) principales :

\begin{enumerate}
  \item \textbf{Stage \texttt{test}} : exécution parallèle de tous les types de tests.
  \item \textbf{Stage \texttt{report}} : génération du rapport Allure consolidé.
\end{enumerate}

\begin{lstlisting}[language=bash, caption={Structure du pipeline CI -- .gitlab-ci-tests.yml (extrait)}, label={lst:gitlab-ci}]
# Template de base pour les tests Playwright
.test-playwright-base:
  image: mcr.microsoft.com/playwright:v1.49.0-noble
  stage: test
  before_script:
    - cd tests/playwright
    - npm ci
    - cp .env.example .env
  artifacts:
    when: always
    paths:
      - tests/playwright/allure-results/
    reports:
      junit: tests/playwright/test-results/results.xml

# Seed des donnees de test
seed-test-data:
  stage: test
  image: node:20
  script:
    - cd tests/seed
    - npm ci
    - node seed-test-data.js

# Tests unitaires backend
test-robot:backend-unit-tests:
  stage: test
  image: maven:3.9-eclipse-temurin-17
  script:
    - cd backend
    - mvn test -pl services/auth-service,
      services/order-service,services/user-service
      --fail-at-end

# Tests API
test-robot:api:
  extends: .test-playwright-base
  script:
    - npx playwright test --project=api-tests
  needs:
    - job: seed-test-data
      optional: true

# Tests E2E Web
test-robot:admin-panel:
  extends: .test-playwright-base
  script:
    - npx playwright test --project=admin-panel

# Tests E2E Mobile
test-robot:customer-app:
  stage: test
  image: cirrusci/flutter:stable
  script:
    - cd frontend/customer_app
    - flutter pub get
    - flutter test integration_test/app_test.dart

# Tests de performance
test-robot:performance-load:
  stage: test
  image:
    name: grafana/k6:latest
    entrypoint: ['']
  script:
    - k6 run tests/performance/k6-load-test.js

# Rapport Allure
test-robot:allure-report:
  stage: report
  image: node:20
  script:
    - npm install -g allure-commandline
    - allure generate tests/playwright/allure-results
      --clean -o allure-report
  needs:
    - job: test-robot:api
    - job: test-robot:admin-panel
\end{lstlisting}

\subsection{Stratégie de déclenchement}

Chaque job est configuré avec des règles de déclenchement spécifiques :

\begin{itemize}
  \item \textbf{Tests unitaires backend} : déclenchés sur les \textit{merge requests}
        modifiant le dossier \texttt{backend/} et sur les branches \texttt{develop}
        et \texttt{staging}.
  \item \textbf{Tests API} : déclenchés sur toutes les \textit{merge requests} et
        les branches \texttt{develop}, \texttt{staging} et les tags de staging.
  \item \textbf{Tests E2E} : déclenchés lorsque les fichiers du frontend ou du
        backend concerné sont modifiés.
  \item \textbf{Tests de performance} : exécutés uniquement sur la branche
        \texttt{staging} et les tags de déploiement, car ils nécessitent un
        environnement proche de la production.
  \item \textbf{Tests mobile} : exécutés avec \texttt{allow\_failure: true} car
        ils dépendent de la disponibilité d'un émulateur dans l'environnement CI.
\end{itemize}

\subsection{Gestion des dépendances entre jobs}

Le mot-clé \texttt{needs} de GitLab CI est utilisé pour définir les dépendances :
les tests API et E2E dépendent du job \texttt{seed-test-data} (avec
\texttt{optional: true} pour ne pas bloquer si le seed échoue), et le rapport Allure
dépend de tous les jobs de test pour agréger leurs résultats.

% =============================================================================
\section{Reporting avec Allure}
\label{sec:reporting-allure}
% =============================================================================

Allure est un outil de reporting qui génère des rapports de tests interactifs à partir
des résultats bruts. Son intégration dans notre projet permet de visualiser les résultats
de manière structurée et de faciliter le diagnostic des échecs.

\subsection{Configuration dans Playwright}

Le listing~\ref{lst:playwright-config} montre la configuration du reporter Allure
dans le fichier \texttt{playwright.config.ts}.

\begin{lstlisting}[language=Java, caption={Configuration Allure dans playwright.config.ts}, label={lst:playwright-config}]
import { defineConfig, devices } from '@playwright/test';

export default defineConfig({
  testDir: '.',
  fullyParallel: true,
  forbidOnly: !!process.env.CI,
  retries: process.env.CI ? 2 : 0,
  workers: process.env.CI ? 2 : undefined,
  timeout: 60_000,

  reporter: [
    ['html', { open: 'never' }],
    ['allure-playwright'],
    ['list'],
  ],

  use: {
    trace: 'on-first-retry',
    screenshot: 'only-on-failure',
    video: 'on-first-retry',
    actionTimeout: 15_000,
    navigationTimeout: 30_000,
  },

  projects: [
    {
      name: 'admin-panel',
      testDir: './e2e/admin-panel',
      use: {
        ...devices['Desktop Chrome'],
        baseURL: ADMIN_BASE_URL,
        storageState: 'fixtures/.auth/admin.json',
      },
      dependencies: ['admin-auth-setup'],
    },
    {
      name: 'api-tests',
      testDir: './api',
      use: { baseURL: API_GATEWAY_URL },
    },
  ],
});
\end{lstlisting}

\subsection{Fonctionnalités du rapport Allure}

Le rapport Allure généré offre les fonctionnalités suivantes :

\begin{itemize}
  \item \textbf{Vue d'ensemble} : tableau de bord avec le nombre de tests
        réussis, échoués, ignorés et cassés.
  \item \textbf{Suites de tests} : organisation hiérarchique par fichier et
        par \texttt{describe} bloc.
  \item \textbf{Chronologie} : visualisation temporelle de l'exécution des tests,
        permettant d'identifier les goulots d'étranglement.
  \item \textbf{Catégories d'erreurs} : classification automatique des échecs
        par type (défaut produit, défaut de test, infrastructure).
  \item \textbf{Pièces jointes} : captures d'écran, traces Playwright et vidéos
        automatiquement attachées aux tests échoués grâce à la configuration
        \texttt{screenshot: 'only-on-failure'} et \texttt{trace: 'on-first-retry'}.
  \item \textbf{Historique} : comparaison avec les exécutions précédentes pour
        suivre l'évolution de la stabilité des tests.
\end{itemize}

La génération du rapport est automatisée dans le pipeline CI via le job
\texttt{test-robot:allure-report}, qui agrège les résultats de tous les jobs de test
et produit un rapport hébergé comme artefact GitLab accessible pendant 30 jours.

% =============================================================================
\section*{Conclusion}
% =============================================================================

Dans ce chapitre, nous avons présenté la réalisation concrète du robot de tests
automatisés pour la plateforme SpeedLine. Au total, notre implémentation comprend :

\begin{itemize}
  \item \textbf{419 tests unitaires} backend couvrant les 11 microservices Java.
  \item \textbf{34 tests unitaires} frontend pour les applications Angular.
  \item \textbf{104 tests API} vérifiant les points d'entrée REST à travers 17
        fichiers de spécification.
  \item \textbf{Tests E2E} web (Playwright) et mobiles (Flutter + Patrol) couvrant
        les parcours utilisateur critiques.
  \item \textbf{3 types de tests de performance} (charge, stress, pic) avec k6.
  \item \textbf{16 tests de sécurité} couvrant les injections SQL, le XSS, le
        contrôle d'accès et la limitation de débit.
  \item Un \textbf{pipeline CI/CD} GitLab automatisé avec reporting Allure.
\end{itemize}

L'ensemble de ces tests s'intègre dans une architecture cohérente respectant la
pyramide de tests définie au chapitre précédent. Le chapitre suivant présentera
les résultats détaillés de l'exécution de ces tests et leur validation.
